\section{Математические модели в экономике}
\label{sec:models-in-economics}

\textbf{Математическая модель в экономике} --- это описание экономического 
процесса или явления с помощью математических зависимостей, уравнений, 
неравенств и функций. Такие модели позволяют формализовать реальные 
экономические процессы, определить влияние различных факторов и найти 
оптимальные решения.

С помощью математических моделей решаются следующие основные задачи:

\begin{itemize}
	\item \textbf{прогнозирование} --- например, прогнозирование спроса, цен и прибыли;
	\item \textbf{оптимизация} --- поиск наилучшего варианта использования ограниченных ресурсов;
	\item \textbf{анализ взаимосвязей} между экономическими показателями;
	\item \textbf{планирование производства и распределения ресурсов};
	\item \textbf{оценка рисков} и различных вариантов развития экономики.
\end{itemize}

Можно выделить несколько распространённых типов экономико-математических задач.

\subsection{Задача оптимизации производства}

Предприятие производит несколько видов продукции. Необходимо определить 
объёмы производства, чтобы получить максимальную прибыль при ограниченных 
ресурсах.

Математическая постановка:

\[
\max Z = \sum_{i=1}^{n} p_i x_i
\]

при ограничениях

\[
\sum_{i=1}^{n} a_{ji}x_i \leq b_j,
\qquad j=1,\ldots,m,
\]

\[
x_i \geq 0.
\]

Здесь $x_i$ --- объём производства $i$-го товара, $p_i$ --- прибыль от 
единицы продукции, $a_{ji}$ --- количество ресурса $j$, необходимое для 
производства единицы товара, $b_j$ --- имеющийся запас ресурса.

\textbf{Смысл модели:} определить, сколько продукции каждого вида следует 
производить, чтобы максимально использовать имеющиеся ресурсы и получить 
наибольшую прибыль.

\subsection{Транспортная задача}

Предположим, имеется несколько поставщиков и несколько потребителей. 
Требуется определить объёмы перевозок так, чтобы все потребности были 
удовлетворены, а общие транспортные затраты были минимальными.

Математическая постановка:

\[
\min Z =
\sum_{i=1}^{m}\sum_{j=1}^{n} c_{ij}x_{ij}
\]

при условиях

\[
\sum_{j=1}^{n}x_{ij}=a_i,
\]

\[
\sum_{i=1}^{m}x_{ij}=b_j,
\]

\[
x_{ij}\geq0.
\]

Здесь $x_{ij}$ --- количество груза, перевозимого от поставщика $i$ 
к потребителю $j$, $c_{ij}$ --- стоимость перевозки единицы груза, 
$a_i$ --- запас поставщика, $b_j$ --- потребность потребителя.

\textbf{Смысл модели:} найти наиболее дешёвую схему распределения 
и перевозки товаров.

\subsection{Модель спроса и предложения}

Математически можно описать зависимость спроса и предложения от цены товара:

\[
Q_d=a-bP,
\]

\[
Q_s=c+dP,
\]

где $P$ --- цена товара, $Q_d$ --- величина спроса, 
$Q_s$ --- величина предложения.

Рыночное равновесие определяется условием

\[
Q_d=Q_s.
\]

\textbf{Смысл модели:} определить равновесную цену и объём продаж, 
при которых количество товара, которое хотят купить потребители, 
совпадает с количеством товара, которое готовы предложить производители.

\subsection{Модель экономического роста}

Для описания изменения выпуска продукции во времени может использоваться, 
например, простая модель:

\[
Y(t)=Y_0e^{gt},
\]

где $Y(t)$ --- объём производства в момент времени $t$, 
$Y_0$ --- начальный объём производства, $g$ --- темп экономического роста.

\textbf{Смысл модели:} оценить, как изменяется объём производства 
при заданном темпе роста и спрогнозировать экономические показатели 
на будущее.

\subsection{Стохастические модели}

В отличие от детерминированных моделей, в \textbf{стохастических моделях} 
учитывается случайный характер экономических процессов. Некоторые параметры 
модели рассматриваются как случайные величины, распределение которых можно 
описать с помощью вероятностей.

\textbf{Стохастическая задача оценки европейского опциона}

Одной из классических задач финансовой математики является оценка 
стоимости европейского опциона. В модели Блэка--Шоулза цена базового 
актива рассматривается как случайный процесс, описываемый геометрическим 
броуновским движением:

\[
dS_t = \mu S_t\,dt + \sigma S_t\,dW_t,
\]

где $S_t$ --- цена базового актива в момент времени $t$, 
$\mu$ --- ожидаемый темп роста цены, $\sigma$ --- волатильность, 
$W_t$ --- стандартное броуновское движение.

Рассмотрим европейский колл-опцион со сроком исполнения $T$ 
и ценой исполнения $K$. Его выплата в момент $T$ является случайной 
величиной:

\[
H=(S_T-K)^+.
\]

При постоянных $r$ и $\sigma$ эта задача приводит к формуле 
Блэка--Шоулза для европейского колл-опциона:

\[
V_t =
S_t N(d_1)
-
K e^{-r(T-t)}N(d_2),
\]

где$r$ --- безрисковая процентная ставка,

\[
d_{1,2} =
\frac{
	\ln(S_t/K)+(r \pm \frac{\sigma^2}{2})(T-t)
}{
	\sigma\sqrt{T-t}
},
\]

а $N(\cdot)$ --- функция распределения стандартной нормальной случайной 
величины.

\textbf{Смысл модели:} цена опциона определяется с учётом случайного 
характера будущей цены базового актива. Модель позволяет найти 
теоретическую стоимость опциона на текущий момент, исходя из его 
текущей цены, цены исполнения, срока до погашения, волатильности 
и безрисковой процентной ставки.

\newpage